\documentclass[12pt, a4paper]{article}

% ── Shared preamble (src/lecturenotes.sty) ────────────────────────────────────
%   Packages, colors, boxes, \sol/\soleq and math macros all live in
%   src/*.sty. \input@path points LaTeX two folders up to find them.
\makeatletter
\def\input@path{{../../}}
\makeatother
\usepackage{lecturenotes}

% ── Local macros ──────────────────────────────────────────────────────────────
\providecommand{\bh}{\mathbf{h}}
\newcommand{\tok}[1]{\textsf{#1}}   % a word / token as the model sees it

% ── Local TikZ styles ─────────────────────────────────────────────────────────
\usetikzlibrary{decorations.pathreplacing}
\tikzset{
  layerlabel/.style={font=\small\itshape, gray},
  brace/.style={decorate, decoration={brace, amplitude=4pt}, thick, gray},
  word/.style={draw=gray!70, fill=white, rounded corners=2pt,
      minimum height=0.65cm, inner xsep=4pt, font=\small\sffamily},
  cell/.style={draw=mainblue, fill=boxbg, thick, rounded corners=3pt,
      minimum size=0.9cm, font=\small},
  opbox/.style={draw=orange!70!black, fill=notebg, thick, rounded corners=4pt,
      minimum height=1.1cm, text width=3.2cm, font=\small, align=center},
  era/.style={circle, fill=mainblue, inner sep=2pt},
}

% Bullet lists: one short sentence per item, little vertical space
\setlist[itemize]{nosep}

% ── Document ──────────────────────────────────────────────────────────────────
\begin{document}

\begin{center}
  {\LARGE \textbf{Section 5.1: From RNN to LLM}} \\[0.3em]
  {\large\itshape A Roadmap of Language Models} \\[0.5em]
  {\large AI-IM Class --- Lecture Notes}%
  \ifstudentmode
    \quad {\normalsize\color{gray}[Student Copy]}%
  \fi
\end{center}

\vspace{1em}
\hrule
\vspace{1.5em}

\begin{itemize}
  \item In Chapter 4 the input was an \textbf{image}; now it is
        \textbf{text}.
  \item This is a bird's-eye view of 30 years of research, up to ChatGPT
        and Claude.
  \item One question runs through it all:
        \emph{how can a computer predict the next word well?}
  \item Each step fixes a weakness of the step before it.
\end{itemize}

\paragraph{Roadmap.}
\begin{itemize}
  \item First we define the task: \textbf{next-word prediction}.
  \item Then we follow the models: N-grams $\to$ embeddings $\to$ RNN $\to$
        Transformer.
  \item We see what \textbf{scale} did to Transformers (GPT-3).
  \item We see how a text predictor became an \textbf{assistant} (ChatGPT).
  \item Finally, we look at where the field is now.
\end{itemize}

\begin{figure}[H]
  \centering
  \begin{tikzpicture}[x=1.45cm, >=Stealth]
    \draw[->, very thick, gray] (-0.3,0) -- (10.4,0);
    \foreach \x/\yr/\lbl/\pos in {
        0/1990s/N-grams/above,
        1.2/1997/LSTM/below,
        2.4/2013/word2vec/above,
        3.6/2014/{seq2seq,\\attention}/below,
        4.8/2017/Transformer/above,
        6.0/2018/{BERT, GPT-1}/below,
        7.2/2020/GPT-3/above,
        8.4/2022/ChatGPT/below,
        9.6/2024+/{reasoning,\\agents}/above} {
        \node[era] at (\x,0) {};
        \ifnum\pdfstrcmp{\pos}{above}=0
          \node[above=3pt, font=\scriptsize\bfseries] at (\x,0) {\yr};
          \node[above=16pt, font=\scriptsize, align=center] at (\x,0) {\lbl};
        \else
          \node[below=3pt, font=\scriptsize\bfseries] at (\x,0) {\yr};
          \node[below=16pt, font=\scriptsize, align=center, anchor=north]
          at (\x,0) {\lbl};
        \fi
      }
  \end{tikzpicture}
  \caption{The roadmap of this section.}
  \label{fig:timeline}
\end{figure}

% ══════════════════════════════════════════════════════════════════════════════
\section{The Task: Predict the Next Word}
\label{sec:task}

\paragraph{NLP and LLMs.}
\begin{itemize}
  \item \textbf{Natural language processing} (NLP) is the part of AI that
        works with human language.
  \item Its tasks include translation, question answering, summarisation and
        speech-to-text.
  \item A \textbf{large language model} (LLM) is one model that can do almost
        all of them.
\end{itemize}

\paragraph{What is the ``label'' for text?}
\begin{itemize}
  \item An image classifier predicts the right label for an image.
  \item A text model is given the words so far and
        \textbf{predicts the next word}.
\end{itemize}

\begin{figure}[H]
  \centering
  \begin{tikzpicture}[>=Stealth]
    \node[word] (a) at (0,0) {The};
    \node[word, right=2pt of a] (b) {sky};
    \node[word, right=2pt of b] (c) {is};
    \node[word, right=2pt of c, draw=orange!80!black, thick, dashed,
      minimum width=1.2cm] (d) {?};
    \node[opbox, right=1.0cm of d, text width=2.2cm] (lm) {language\\model};
    \draw[->, thick] (d) -- (lm);
    \begin{scope}[shift={($(lm.east)+(2.0,1.1)$)}]
      \foreach \w/\p [count=\i] in {blue/3.5, clear/2.2, the/1.6,
          falling/0.9, red/0.15} {
          \node[anchor=east, font=\small\sffamily] at (0,-\i*0.45) {\w};
          \fill[mainblue!70] (0.1,-\i*0.45-0.15) rectangle ++(\p,0.3);
        }
      \node[font=\scriptsize, gray, anchor=west] at (0.1,-2.75)
      {probability of each next word};
    \end{scope}
    \draw[->, thick] (lm.east) -- ++(0.6,0);
  \end{tikzpicture}
  \caption{The model gives a probability to every word in its vocabulary
    (made-up numbers).}
  \label{fig:nextword}
\end{figure}

\begin{resultbox}
  A \textbf{language model} computes
  \[
    P(\text{next word} \mid \text{words so far}).
  \]
  To generate text, it picks a likely word, appends it, and repeats.
\end{resultbox}

\begin{itemize}
  \item This is just \textbf{classification}: the classes are the words, with
        a softmax and cross-entropy.
  \item Every model in this section has this same output.
  \item The models differ only in \emph{how they read the words so far}.
\end{itemize}

\paragraph{Tokens.}
\begin{itemize}
  \item Models see \textbf{tokens}, not whole words.
  \item A frequent word is one token; a rare word is split into pieces, e.g.\
        \tok{unbelievable} $\to$ \tok{un} + \tok{believ} + \tok{able}.
  \item A modern vocabulary has about 50\,000 to 200\,000 tokens.
  \item In these notes, ``word'' and ``token'' mean the same thing.
\end{itemize}

\begin{notebox}
  \textbf{Try it.} \url{https://tiktokenizer.vercel.app}
  \begin{itemize}
    \item Paste one English and one Thai sentence.
    \item Compare the token counts: more tokens means more cost.
  \end{itemize}
\end{notebox}

% ══════════════════════════════════════════════════════════════════════════════
\section{N-grams: Counting}
\label{sec:ngram}

\begin{itemize}
  \item The oldest approach is to read a lot of text and \textbf{count}.
\end{itemize}
\[
  P(\tok{blue} \mid \tok{the sky is})
  = \frac{\text{count}(\tok{the sky is blue})}{\text{count}(\tok{the sky is})}
\]
\begin{itemize}
  \item An \textbf{N-gram} model looks only at the last $N - 1$ words.
  \item The formula above is a 4-gram model.
\end{itemize}

\begin{stepbox}{Exercise 1 \quad An N-gram by hand}
  In a corpus, \tok{the sky is} appears 200 times, followed by \tok{blue}
  120 times, \tok{clear} 50 times and \tok{grey} 30 times.
  \begin{enumerate}[label=(\alph*), nosep]
    \item $P(\tok{blue} \mid \tok{the sky is}) = \sol[3cm]{120/200 = 0.6}$.
    \item $P(\tok{red} \mid \tok{the sky is}) = \sol[3cm]{0}$, because
          $\sol[6cm]{\text{it never appears in the corpus}}$.
  \end{enumerate}
\end{stepbox}

\paragraph{Weaknesses.}
\begin{itemize}
  \item \textbf{Too specific}: \tok{the sky was} and \tok{the sky is} are
        counted as unrelated.
  \item \textbf{Short memory}: a 4-gram forgets everything more than three
        words back.
  \item A larger $N$ does not help, because most long sequences never appear
        in the data.
\end{itemize}

% ══════════════════════════════════════════════════════════════════════════════
\section{Word Embeddings: Words as Vectors}
\label{sec:embed}

\begin{itemize}
  \item To fix ``too specific'', we turn each word into a vector, called an
        \textbf{embedding}.
  \item Words with similar meaning get similar vectors.
  \item \textbf{Word2vec} (2013) learns them by predicting a word from its
        neighbours.
  \item Words used in similar contexts end up close together, e.g.\
        \tok{is}/\tok{was} or \tok{cat}/\tok{dog}.
\end{itemize}

\begin{figure}[H]
  \centering
  \begin{tikzpicture}[x=1.1cm, y=1.1cm, >=Stealth]
    \draw[->, gray] (0,0) -- (6,0);
    \draw[->, gray] (0,0) -- (0,3.6);
    \foreach \x/\y/\w in {1/0.8/man, 2/2.6/woman, 4/1.0/king, 5/2.8/queen} {
        \fill[mainblue] (\x,\y) circle (2.5pt);
        \node[right=2pt, font=\small\sffamily] at (\x,\y) {\w};
      }
    \draw[->, thick, orange!80!black] (1,0.8) -- (1.95,2.5);
    \draw[->, thick, orange!80!black] (4,1.0) -- (4.95,2.7);
    \node[font=\scriptsize, orange!80!black, align=center] at (2.9,1.5)
    {same direction:\\``male $\to$ female''};
    \node[layerlabel] at (3,-0.45) {a 2-D picture of embedding space
      (real embeddings have hundreds of dimensions)};
  \end{tikzpicture}
  \caption{Directions in embedding space carry meaning.}
  \label{fig:embed}
\end{figure}

\begin{stepbox}{Exercise 2 \quad Vector arithmetic on words}
  Using Figure~\ref{fig:embed}:
  $\;\tok{king} - \tok{man} + \tok{woman} \approx \sol[3cm]{\tok{queen}}$.
\end{stepbox}

\begin{itemize}
  \item Embeddings are still the \textbf{first layer} of every LLM today.
\end{itemize}

\paragraph{Weakness: no context.}
\begin{itemize}
  \item A word gets the same vector in every sentence.
  \item \tok{bank} looks the same in \emph{river bank} and in
        \emph{bank account}.
  \item We need a model that reads the whole sentence.
\end{itemize}

\begin{notebox}
  \textbf{Try it.} \url{https://projector.tensorflow.org}
  \begin{itemize}
    \item Search for a word and see its nearest neighbours in 3-D.
  \end{itemize}
\end{notebox}

% ══════════════════════════════════════════════════════════════════════════════
\section{Recurrent Neural Networks: Memory}
\label{sec:rnn}

\begin{itemize}
  \item A \textbf{recurrent neural network} (RNN) reads one word at a time.
  \item It keeps a \textbf{hidden state} $\bh_t$ that summarises everything
        read so far.
  \item At each step it combines the old memory with the new word.
\end{itemize}

\begin{resultbox}
  \[
    \bh_t = \tanh\!\big(\bW_{h}\,\bh_{t-1} + \bW_{x}\,\bx_t + \bb\big)
  \]
  \begin{itemize}
    \item $\bx_t$ is the embedding of word $t$.
    \item The \emph{same} $\bW_h, \bW_x, \bb$ are used at every step.
    \item A softmax on top of $\bh_t$ predicts the next word.
  \end{itemize}
\end{resultbox}

\begin{figure}[H]
  \centering
  \begin{tikzpicture}[>=Stealth, node distance=1.6cm]
    \node[cell] (h1) {$\bh_1$};
    \node[cell, right=of h1] (h2) {$\bh_2$};
    \node[cell, right=of h2] (h3) {$\bh_3$};
    \node[left=1.0cm of h1, font=\small] (h0) {$\bh_0$};
    \foreach \i/\w in {1/The, 2/sky, 3/is} {
        \node[word, below=0.7cm of h\i] (x\i) {\w};
        \draw[->, thick] (x\i) -- (h\i);
      }
    \node[font=\small\sffamily, above=0.6cm of h1] (y1) {sky?};
    \node[font=\small\sffamily, above=0.6cm of h2] (y2) {is?};
    \node[font=\small\sffamily, above=0.6cm of h3, orange!80!black] (y3)
    {\textbf{blue?}};
    \foreach \i in {1,2,3} \draw[->, thick, gray] (h\i) -- (y\i);
    \draw[->, thick, mainblue] (h0) -- (h1);
    \draw[->, thick, mainblue] (h1) -- node[above, font=\scriptsize] {memory} (h2);
    \draw[->, thick, mainblue] (h2) -- node[above, font=\scriptsize] {memory} (h3);
  \end{tikzpicture}
  \caption{An RNN unrolled over three words; every box has the same weights.}
  \label{fig:rnn}
\end{figure}

\begin{patternbox}
  \textbf{Weight sharing again.}
  \begin{itemize}
    \item A CNN reuses one kernel at every \emph{position in space}.
    \item An RNN reuses one set of weights at every \emph{position in time}.
    \item In both, the number of parameters does not depend on the input
          size.
  \end{itemize}
\end{patternbox}

\paragraph{LSTM and GRU.}
\begin{itemize}
  \item A plain RNN forgets quickly.
  \item The \textbf{LSTM} (1997) and \textbf{GRU} (2014) add \emph{gates}
        that learn what to keep and what to forget.
  \item LSTMs were the state of the art from about 2014 to 2017.
  \item In 2016 Google Translate switched to an LSTM
        \textbf{encoder--decoder} (seq2seq).
\end{itemize}

\paragraph{Weaknesses.}
\begin{itemize}
  \item \textbf{Memory bottleneck}: the whole history must fit in one vector
        $\bh_t$.
  \item \textbf{Vanishing gradient}: the gradient is multiplied by $\bW_h$ at
        every step, so it shrinks towards zero.
  \item \textbf{Slow}: $\bh_t$ needs $\bh_{t-1}$, so the GPU cannot work in
        parallel.
\end{itemize}

\begin{notebox}
  \textbf{Why is an RNN slow?}
  \begin{itemize}
    \item To compute $\bh_{100}$, the RNN must first compute
          $\bh_1, \bh_2, \dots, \bh_{99}$, one after another.
    \item These steps cannot run in parallel.
  \end{itemize}
\end{notebox}

\begin{notebox}
  \textbf{Read more.}
  \begin{itemize}
    \item LSTM gates, with diagrams:\newline
          \url{https://colah.github.io/posts/2015-08-Understanding-LSTMs/}
    \item What a character-level RNN can generate:\newline
          \url{https://karpathy.github.io/2015/05/21/rnn-effectiveness/}
  \end{itemize}
\end{notebox}

% ══════════════════════════════════════════════════════════════════════════════
\section{Attention and the Transformer}
\label{sec:transformer}

\paragraph{Attention.}
\begin{itemize}
  \item In 2014, attention let RNN translators ``look back'' at any input
        word.
  \item In 2017, \emph{Attention Is All You Need} dropped the RNN and kept
        only attention.
  \item The result is the \textbf{Transformer}.
\end{itemize}

\paragraph{The idea: self-attention.}
\begin{itemize}
  \item Every word looks at every other word in the sentence.
  \item It decides how relevant each of the other words is.
  \item Its new vector is a \textbf{weighted average} of their vectors.
\end{itemize}

\begin{resultbox}
  \[
    \text{new vector of word } i
    = \sum_{j} a_{ij}\,(\text{vector of word } j),
    \qquad a_{ij} \ge 0,\;\; \sum_j a_{ij} = 1
  \]
  \begin{itemize}
    \item The weight $a_{ij}$ is a softmax of how well words $i$ and $j$
          ``match''.
    \item The network learns what a good match is.
  \end{itemize}
\end{resultbox}

\begin{figure}[H]
  \centering
  \begin{tikzpicture}[>=Stealth]
    \node[word] (w1) at (0,0) {The};
    \foreach \w [count=\i from 2] in {animal, didn't, cross, the, street,
        because, it, was, tired} {
        \pgfmathtruncatemacro{\p}{\i-1}
        \node[word, right=3pt of w\p] (w\i) {\w};
      }
    \draw[->, line width=3pt, orange!80!black, opacity=0.85]
    (w8.north) to[out=110, in=70] (w2.north);
    \draw[->, line width=1.2pt, orange!80!black, opacity=0.5]
    (w8.north) to[out=120, in=60] (w6.north);
    \draw[->, line width=0.6pt, orange!80!black, opacity=0.5]
    (w8.south) to[out=-60, in=-120] (w10.south);
    \node[font=\scriptsize, orange!80!black, below=0.55cm of w4]
    {\tok{it} attends mostly to \tok{animal}};
  \end{tikzpicture}
  \caption{Self-attention for \tok{it}; a thicker arrow means a larger
    weight.}
  \label{fig:attention}
\end{figure}

\paragraph{Why it won.}
\begin{itemize}
  \item \textbf{No bottleneck}: word 1 and word 1000 are linked in one step.
  \item \textbf{Parallel}: all words are processed at once, which suits GPUs.
  \item \textbf{Context}: \tok{bank} gets a different vector near
        \tok{river} than near \tok{account}.
\end{itemize}

\paragraph{Inside.}
\begin{itemize}
  \item A Transformer is a stack of identical \textbf{blocks}.
  \item Each block is a self-attention layer followed by a small MLP.
  \item Residual connections, as in ResNet, link the blocks.
  \item A \textbf{positional encoding} tells the model where each word is.
\end{itemize}

\begin{notebox}
  \textbf{RNN versus Transformer.}
  \begin{center}
    \renewcommand{\arraystretch}{1.4}
    \begin{tabular}{lll}
      \toprule
                                    & RNN           & Transformer \\
      \midrule
      Words processed               & one at a time & all at once \\
      Path from word 1 to word 1000 & 999 steps     & 1 step      \\
      Uses the GPU well?            & no            & yes         \\
      \bottomrule
    \end{tabular}
  \end{center}
\end{notebox}

\begin{notebox}
  \textbf{For the curious.}
  The full formula is $\operatorname{softmax}\!\big(QK^{\top}/\sqrt{d}\big)\,V$.
  \begin{itemize}
    \item The \emph{query} asks: what am I looking for?
    \item The \emph{key} says: what do I contain?
    \item The \emph{value} is what gets passed on.
    \item The weight $a_{ij}$ comes from query$_i \cdot$ key$_j$.
  \end{itemize}
\end{notebox}

\paragraph{Three families.}
\begin{itemize}
  \item The original Transformer has an \textbf{encoder} (reads) and a
        \textbf{decoder} (writes).
  \item Later models kept one or both halves.
\end{itemize}

\begin{center}
  \renewcommand{\arraystretch}{1.35}
  \begin{tabular}{>{\raggedright}p{2.7cm}>{\raggedright}p{3.0cm}>{\raggedright}p{3.4cm}>{\raggedright\arraybackslash}p{4.0cm}}
    \toprule
    Family           & Example (year)            & Sees                           & Good at                               \\
    \midrule
    Encoder--decoder & original, T5 (2017, 2019) & the input, then its own output & translation, summarisation            \\
    Encoder only     & BERT (2018)               & the whole sentence             & understanding: classification, search \\
    Decoder only     & GPT (2018)                & only the words so far          & \textbf{generating text}              \\
    \bottomrule
  \end{tabular}
\end{center}

\begin{itemize}
  \item ChatGPT, Claude, Gemini and Llama are all \textbf{decoder-only}.
  \item They do exactly the next-word prediction of Section~\ref{sec:task}.
  \item GPT stands for \emph{Generative Pre-trained Transformer}.
\end{itemize}

\begin{notebox}
  \textbf{Try it.}
  \begin{itemize}
    \item Watch a small GPT run step by step in your browser:\newline
          \url{https://poloclub.github.io/transformer-explainer/}
    \item Walk through a whole LLM in 3-D: \url{https://bbycroft.net/llm}
    \item See the attention weights of a real model:\newline
          \url{https://huggingface.co/spaces/exbert-project/exbert}
  \end{itemize}
\end{notebox}

% ══════════════════════════════════════════════════════════════════════════════
\section{Scale: From GPT-1 to GPT-3}
\label{sec:scale}

\begin{itemize}
  \item Once the architecture was fixed, progress came from making it
        \textbf{bigger}.
\end{itemize}

\begin{center}
  \renewcommand{\arraystretch}{1.35}
  \begin{tabular}{lll}
    \toprule
    Model & Year & Parameters  \\
    \midrule
    GPT-1 & 2018 & 0.1 billion \\
    GPT-2 & 2019 & 1.5 billion \\
    GPT-3 & 2020 & 175 billion \\
    \bottomrule
  \end{tabular}
\end{center}

\paragraph{Scaling laws.}
\begin{itemize}
  \item The error falls smoothly and \emph{predictably} as three things grow
        together:
        \begin{itemize}
          \item the number of parameters,
          \item the amount of training text,
          \item the computing power.
        \end{itemize}
  \item This turned AI progress into an engineering plan: spend more, get a
        better model.
\end{itemize}

\paragraph{New abilities.}
\begin{itemize}
  \item Predicting the next word on the internet also teaches grammar, facts,
        translation and some reasoning.
  \item That is because all of these help to predict the next word.
  \item GPT-3 showed \textbf{in-context learning}: it learns from examples in
        the prompt, without retraining.
\end{itemize}
\begin{center}
  \tt\small
  \begin{tabular}{l}
    sea otter => loutre de mer \\
    cheese => fromage          \\
    peppermint => \textcolor{orange!80!black}{\textit{(model writes:) menthe poivr\'ee}}
  \end{tabular}
\end{center}
\begin{itemize}
  \item This made \textbf{prompting} a new way to program a model.
\end{itemize}

\paragraph{Weakness: a base model is not an assistant.}
\begin{itemize}
  \item A pre-trained model only \emph{continues} text.
  \item Asked ``What is the capital of Thailand?'', it may just write more
        quiz questions.
\end{itemize}

% ══════════════════════════════════════════════════════════════════════════════
\section{From Text Predictor to Assistant}
\label{sec:assistant}

\begin{itemize}
  \item ChatGPT (November 2022) added two training stages after
        pre-training.
  \item Every chat model since then follows the same recipe.
\end{itemize}

\begin{figure}[H]
  \centering
  \begin{tikzpicture}[>=Stealth, node distance=0.6cm]
    \tikzset{stage/.style={opbox, text width=3.6cm, minimum height=2.2cm}}
    \node[stage] (pt) {\textbf{1. Pre-training}\\[4pt]\scriptsize
      next-word prediction on\\trillions of tokens of text};
    \node[stage, right=of pt] (sft) {\textbf{2. Instruction}\\
      \textbf{tuning}\\[4pt]\scriptsize learn from examples\\of good answers};
    \node[stage, right=of sft] (rl) {\textbf{3. Learning}\\
      \textbf{from feedback}\\[4pt]\scriptsize reward answers\\people
      prefer (RLHF)};
    \draw[->, thick] (pt) -- (sft);
    \draw[->, thick] (sft) -- (rl);
    \node[layerlabel, below=0.25cm of pt] {knowledge};
    \node[layerlabel, below=0.25cm of sft] {format: answer, don't continue};
    \node[layerlabel, below=0.25cm of rl] {helpful, honest, safe};
  \end{tikzpicture}
  \caption{The three training stages; stage 1 uses most of the data and
    compute.}
  \label{fig:stages}
\end{figure}

\begin{itemize}
  \item \textbf{Pre-training} gives the model language and knowledge about
        the world.
  \item \textbf{Instruction tuning} (SFT) trains it on example conversations
        written by people.
  \item \textbf{RLHF} (reinforcement learning from human feedback):
        \begin{itemize}
          \item The model writes several answers, and people rank them.
          \item The model is then trained towards the higher-ranked answers.
          \item A variant uses AI feedback guided by written principles
                (\emph{Constitutional AI}, Anthropic).
        \end{itemize}
\end{itemize}

\begin{stepbox}{Exercise 3 \quad Which stage?}
  Name the training stage that is mainly responsible for each behaviour.
  \begin{enumerate}[label=(\alph*), nosep]
    \item The model knows that Bangkok is the capital of Thailand:
          $\sol[4cm]{\text{pre-training}}$.
    \item It answers a question instead of writing more questions:
          $\sol[4cm]{\text{instruction tuning}}$.
    \item Of two correct answers, the model prefers the clearer and more
          polite one: $\sol[4cm]{\text{RLHF}}$.
  \end{enumerate}
\end{stepbox}

% ══════════════════════════════════════════════════════════════════════════════
\section{Where the Field Is Now}
\label{sec:now}

\begin{itemize}
  \item The core is unchanged: a decoder-only Transformer trained in three
        stages.
\end{itemize}

\paragraph{Main directions.}
\begin{itemize}
  \item \textbf{Multimodal}: images, audio and documents are turned into
        tokens too.
        \begin{itemize}
          \item A model can read a photo of a machine, a scanned form or a
                chart.
        \end{itemize}
  \item \textbf{Long context}: from about 2\,000 tokens to hundreds of
        thousands or millions.
        \begin{itemize}
          \item A model can read a whole manual or code base at once.
        \end{itemize}
  \item \textbf{Open weights}: anyone can download and run these models.
        \begin{itemize}
          \item Examples are Llama, Mistral, Qwen and DeepSeek, from laptop
                size to frontier size.
        \end{itemize}
  \item \textbf{Efficiency}: large models are made cheaper to run.
        \begin{itemize}
          \item \emph{Mixture of experts} uses only part of the network for
                each token.
          \item \emph{Distillation} trains small models to copy large ones.
        \end{itemize}
  \item \textbf{Reasoning models} (2024 onwards):
        \begin{itemize}
          \item They are trained to ``think'' step by step before answering.
          \item Spending more compute at answer time (\emph{test-time
                compute}) gives better answers.
          \item The biggest gains are in maths, science and programming.
        \end{itemize}
  \item \textbf{Tools and agents}:
        \begin{itemize}
          \item A model can search the web, run code, query databases and
                edit files.
          \item An \textbf{agent} repeats think $\to$ act $\to$ observe until
                the task is done.
          \item Coding agents such as Claude Code handle multi-step
                programming tasks.
        \end{itemize}
\end{itemize}

\begin{warnbox}
  \textbf{Limitations that remain.}
  \begin{itemize}
    \item \textbf{Hallucination}: the model is trained for \emph{likely}
          text, not \emph{true} text, so check its outputs.
    \item \textbf{Knowledge cutoff}: the model knows nothing after its
          training data ends.
          \begin{itemize}
            \item The fix is a search tool, or supplying documents
                  (\emph{RAG}).
          \end{itemize}
    \item \textbf{Cost}: training a frontier model costs hundreds of millions
          of dollars.
  \end{itemize}
\end{warnbox}

% ══════════════════════════════════════════════════════════════════════════════
\section{Summary}
\label{sec:summary}

\begin{center}
  \renewcommand{\arraystretch}{1.4}
  \begin{tabular}{>{\raggedright}p{2.6cm}>{\raggedright}p{5.4cm}>{\raggedright\arraybackslash}p{5.4cm}}
    \toprule
    Step                    & Key idea                                   & Weakness that led to the next step \\
    \midrule
    N-grams                 & count word sequences                       & too specific; short memory         \\
    Embeddings (2013)       & words as vectors; similar meaning is close & one vector per word, no context    \\
    RNN / LSTM (1997--2014) & hidden state carries memory                & vanishing gradient; not parallel   \\
    Transformer (2017)      & self-attention: every word sees every word & ---                                \\
    Scale (2018--2020)      & bigger model, more data, more compute      & only continues text                \\
    Assistant (2022)        & instruction tuning + RLHF                  & hallucination, cutoff, cost        \\
    Now (2023--)            & multimodal, reasoning, tools and agents    &                                    \\
    \bottomrule
  \end{tabular}
\end{center}

\begin{itemize}
  \item Every model solves the same task: classify the next word with a
        softmax.
  \item The models differ in how they read the context: by counting, by a
        memory vector, or by attention.
  \item Attention gives direct links and runs in parallel, so huge models
        became trainable.
  \item Scale gave broad ability; SFT and RLHF turned the model into an
        assistant.
\end{itemize}

\paragraph{Resources.}
\begin{itemize}
  \item Tokenizer playground: \url{https://tiktokenizer.vercel.app}
  \item Embedding projector: \url{https://projector.tensorflow.org}
  \item Understanding LSTMs:\newline
        \url{https://colah.github.io/posts/2015-08-Understanding-LSTMs/}
  \item Transformer Explainer:
        \url{https://poloclub.github.io/transformer-explainer/}
  \item LLM visualisation: \url{https://bbycroft.net/llm}
  \item Attention visualised:
        \url{https://huggingface.co/spaces/exbert-project/exbert}
  \item A full free NLP course: \url{https://www.nlpdemystified.org/course}
\end{itemize}

\end{document}
